\documentclass[a4paper]{article}

\usepackage[fontset=fandol]{ctex}
\usepackage{amsmath, amssymb}
\usepackage{graphicx}
\usepackage[margin=2.5cm]{geometry}
\usepackage[most]{tcolorbox}
\usepackage{etoolbox}
\usepackage{listings}
\usepackage{hyperref}
\usepackage{booktabs}
\usepackage{subcaption}
\usepackage{float}
\usepackage{tikz}
\IfFileExists{pgfplots.sty}{
    \usepackage{pgfplots}
    \pgfplotsset{compat=1.18}
}{}

\newtcolorbox{knowledgebox}[1]{
    enhanced, colback=blue!5!white, colframe=blue!75!black, colbacktitle=blue!75!black,
    coltitle=white, fonttitle=\bfseries, title=#1, attach boxed title to top left={yshift=-2mm, xshift=2mm},
    boxrule=1pt, sharp corners
}

\newtcolorbox{importantbox}[1]{
    enhanced, colback=yellow!10!white, colframe=yellow!80!black, colbacktitle=yellow!80!black,
    coltitle=black, fonttitle=\bfseries, title=#1, sharp corners
}

\newtcolorbox{warningbox}[1]{
    enhanced, colback=red!5!white, colframe=red!75!black, colbacktitle=red!75!black,
    coltitle=white, fonttitle=\bfseries, title=#1, sharp corners
}

\lstset{
    language=Python,
    basicstyle=\ttfamily\small,
    keywordstyle=\color{blue},
    stringstyle=\color{red!60!black},
    commentstyle=\color{green!60!black},
    breaklines=true,
    frame=single,
    numbers=left,
    numberstyle=\tiny\color{gray},
    captionpos=b,
    extendedchars=false
}

\newcommand{\notetitle}{CS231n Lecture 7: Recurrent Neural Networks}
\newcommand{\noteauthors}{Codex Harness}
\newcommand{\notedate}{2026-04-15}
\newcommand{\videochannel}{Stanford Online / CS231n}
\newcommand{\videopublishdate}{2025-04-22}
\newcommand{\videoduration}{01:11:43}
\newcommand{\videourl}{https://www.youtube.com/watch?v=kG2lAPBF7zA}
\newcommand{\videocoverpath}{cover.jpg}

\begin{document}

\begin{titlepage}
\centering
{\Large 课程笔记\par}
\vspace{1.2cm}
{\huge\bfseries \notetitle\par}
\vspace{0.8cm}
{\large \noteauthors\par}
\vspace{0.3cm}
{\large \notedate\par}
\vspace{1.2cm}

\includegraphics[width=0.82\textwidth,height=0.45\textheight,keepaspectratio]{\videocoverpath}\par

\vfill
\begin{tcolorbox}[width=0.9\textwidth, colback=black!2!white, colframe=black!60, sharp corners]
\textbf{视频作者/频道}：\videochannel\par
\textbf{发布日期}：\videopublishdate\par
\textbf{视频时长}：\videoduration\par
\textbf{视频链接}：\href{\videourl}{\nolinkurl{\videourl}}
\end{tcolorbox}
\end{titlepage}

\tableofcontents
\newpage

\section{为什么序列建模需要新的网络结构}
% source_units: cu_07_02
在前几讲中，视觉模型大多假设输入是固定维度的向量、网格或特征图；网络读入一次输入，输出一次预测。但本讲一开始就明确指出，计算机视觉中的很多任务天然是序列（sequence）问题，而不是单次静态映射。典型例子包括：

\begin{itemize}
\item 图像描述（image captioning）：输入一张图像，输出一个单词序列。
\item 动作预测（action prediction）：输入一串视频帧，输出动作类别。
\item 视频描述（video captioning）：输入视频帧序列，输出一句话。
\item 帧级别视频分类：输入随时间变化的特征流，输出每一步或整体标签。
\end{itemize}

这些任务的共同点是：模型不只关心“当前输入是什么”，还必须记住“之前已经看到了什么”。因此，课程把 \textbf{内部状态（internal state）} 作为序列模型的核心思想。换句话说，序列模型不是把整条序列一次性拍平成一个向量，而是在时间维度上持续维护一个可以被更新的记忆。

\begin{knowledgebox}{本讲主线}
本讲不是简单地给出几个新层，而是沿着一条非常清晰的路径推进：先说明序列任务的需求，再引出 vanilla RNN 的参数共享与 hidden state，随后讨论 BPTT 的训练困难，最后用 LSTM/GRU 这样的门控结构改善长程依赖建模。
\end{knowledgebox}

\subsection{本章小结}
序列任务与静态分类任务的差异，不在于输出形式是否更长，而在于模型必须在不同时间步之间持续传递信息。本讲后续所有结构，都是围绕“怎样保存和更新这份状态”展开的。

\section{Vanilla RNN 的状态、参数共享与计算图}
% source_units: cu_07_03, cu_07_04, cu_07_05
课程首先用“展开后的时间轴”来解释 RNN。图中每个时间步都会接收当前输入 $x_t$，并从前一时刻继承 hidden state $h_{t-1}$，再产生新的 hidden state $h_t$ 与输出 $y_t$。

\begin{figure}[H]
\centering
\includegraphics[width=0.82\textwidth]{pdf_pages/page-21.png}
\caption{将 RNN 沿时间展开后，可以直观看到输入、状态与输出之间的依赖关系。}
\end{figure}

最核心的两条公式是：
\[
h_t = \phi(W_{xh}x_t + W_{hh}h_{t-1} + b_h)
\]
\[
y_t = W_{hy}h_t + b_y
\]

其中：
\begin{itemize}
\item $x_t$：第 $t$ 个时间步的输入向量。
\item $h_t$：第 $t$ 个时间步的隐藏状态（hidden state）。
\item $W_{xh}$：输入到隐藏状态的参数矩阵。
\item $W_{hh}$：前一时刻状态到当前状态的递归参数矩阵。
\item $W_{hy}$：隐藏状态到输出的映射矩阵。
\item $\phi(\cdot)$：逐元素非线性函数，课程里常用 $\tanh$ 或 ReLU。
\end{itemize}

\begin{importantbox}{RNN 与普通前馈网络最关键的区别}
同一组参数 $W_{xh}, W_{hh}, W_{hy}$ 会在所有时间步重复使用。这就是所谓的 \textbf{参数共享（parameter sharing）}。它让模型可以处理任意长度序列，同时把“时间结构”编码进网络里。
\end{importantbox}

为了避免公式看起来过于抽象，课程还给了一个很有代表性的玩具例子：在二值序列里检测“重复的 1”。这个例子的重点不是任务本身，而是让读者思考：\textbf{hidden state 到底应该记什么？} 结论很直接，若想判断当前是否又看到了一个 1，状态里至少要编码“上一个输入是不是 1”。这说明 hidden state 不是任意的缓存，而是任务相关信息的压缩表示。

接着，课程把 RNN 计算图分成几种常见模式：
\begin{itemize}
\item \textbf{many-to-many}：输入与输出都是序列，例如逐帧标注。
\item \textbf{many-to-one}：输入是序列，输出是单标签，例如情感分类或动作识别。
\item \textbf{one-to-many}：输入是单个向量，输出是序列，例如图像描述。
\end{itemize}

这种分类非常重要，因为它直接决定损失函数在哪些时间步上定义、反向传播怎样展开，以及推理时需要保留哪些中间状态。

\subsection{本章小结}
Vanilla RNN 用一条反复复用的递归公式，把“当前输入”和“历史状态”融合成新的状态。参数共享赋予它处理变长序列的能力，而 hidden state 则承担了跨时间步的信息压缩任务。

\section{BPTT、截断训练与梯度问题}
% source_units: cu_07_06, cu_07_07
有了前向递归之后，训练就自然落到 \textbf{Backpropagation Through Time, BPTT}。本质上，它并不是发明了新的反向传播，而是先把 RNN 在时间维上展开成一个很深的网络，再对这个展开图做标准反传。

如果把每个时间步的损失记作 $L_t$，那么对状态的梯度递推可以写成：
\[
\frac{\partial L}{\partial h_t}
=
\frac{\partial L_t}{\partial h_t}
+
\frac{\partial L}{\partial h_{t+1}}
\frac{\partial h_{t+1}}{\partial h_t}
\]

继续展开后，会出现一串 Jacobian 连乘：
\[
\frac{\partial L}{\partial h_t}
\propto
\prod_{k=t+1}^{T}
\frac{\partial h_k}{\partial h_{k-1}}
\]

这也是课程后来反复强调 \textbf{vanishing / exploding gradients} 的根源。如果这串矩阵的有效奇异值长期小于 $1$，梯度会指数级衰减；如果长期大于 $1$，梯度又会爆炸。RNN 越想记住更久远的信息，链就越长，问题就越严重。

\begin{figure}[H]
\centering
\includegraphics[width=0.82\textwidth]{pdf_pages/page-51.png}
\caption{Truncated BPTT 用固定窗口回传梯度，是训练长序列时最重要的工程折中之一。}
\end{figure}

因此，课程给出的第一个工程化应对手段不是“换掉 RNN”，而是 \textbf{truncated BPTT}：前向仍然可以持续滚动更长的序列，但反向只在一个有限窗口内展开。这样做的收益很直接：

\begin{itemize}
\item 显著降低显存与计算开销。
\item 让 mini-batch 训练更容易落地。
\item 把优化问题从“无限长依赖”改成“局部窗口中的近似优化”。
\end{itemize}

代价也同样明确：如果真正重要的依赖跨度超出截断窗口，模型就很难从梯度里得到足够信号。

\begin{warningbox}{不要误解 truncated BPTT}
截断不是理论上更正确的训练方式，而是工程上可承受的近似。它解决的是训练可行性，不是表达能力。
\end{warningbox}

课程后半部分还专门用梯度流图指出：vanilla RNN 即便在前向上能“理论上”记住长信息，反向优化也常常先失败。也就是说，难点不只是模型有没有记忆槽位，而是优化器是否真的能把有用信息写进这些槽位。

\begin{figure}[H]
\centering
\includegraphics[width=0.82\textwidth]{pdf_pages/page-105.png}
\caption{沿时间展开的梯度流展示了 recurrent Jacobian 连乘为何会让长依赖训练变得脆弱。}
\end{figure}

\subsection{本章小结}
BPTT 把 RNN 训练问题变成了一个很深的展开网络训练问题。真正的瓶颈不在“能不能写出递推公式”，而在于梯度沿时间传播时是否还能保持稳定。截断 BPTT 是实用近似，梯度消失/爆炸是核心优化障碍。

\section{门控 RNN: LSTM 与 GRU 的核心思想}
% source_units: cu_07_08, cu_07_09, cu_07_10
既然问题出在“所有信息都挤在同一个 hidden state 里”，课程自然引出更复杂的门控结构。LSTM 的核心改动，是把“短时非线性状态”和“更稳定的 cell state”区分开来，再通过若干门（gates）决定信息如何写入、保留和读出。

一种标准写法如下：
\[
f_t = \sigma(W_f[x_t; h_{t-1}] + b_f), \qquad
i_t = \sigma(W_i[x_t; h_{t-1}] + b_i)
\]
\[
\tilde{c}_t = \tanh(W_c[x_t; h_{t-1}] + b_c), \qquad
o_t = \sigma(W_o[x_t; h_{t-1}] + b_o)
\]
\[
c_t = f_t \odot c_{t-1} + i_t \odot \tilde{c}_t, \qquad
h_t = o_t \odot \tanh(c_t)
\]

其中：
\begin{itemize}
\item $f_t$：forget gate，决定旧记忆保留多少。
\item $i_t$：input gate，决定新信息写入多少。
\item $\tilde{c}_t$：candidate memory，表示候选写入内容。
\item $o_t$：output gate，决定当前暴露多少内部记忆作为输出。
\item $c_t$：cell state，是跨时间传递的“主记忆通道”。
\item $\odot$：逐元素乘法。
\end{itemize}

\begin{figure}[H]
\centering
\includegraphics[width=0.9\textwidth]{pdf_pages/page-64.png}
\caption{LSTM 通过 forget/input/output 三类门，把“保留什么、写入什么、暴露什么”显式分开。}
\end{figure}

从教学角度看，LSTM 真正解决的不是“让梯度永不衰减”，而是给网络提供一条\textbf{更容易保留长期信息}的通路。因为 $c_t$ 中包含了 $f_t \odot c_{t-1}$ 这一相对直接的残差式传播路径，只要 forget gate 不把旧状态完全抹掉，跨时间的信息保存就比 vanilla RNN 容易得多。

课程也特别提醒：\textbf{LSTM 并不神奇。} 它只是让长依赖更可训练，而不是在数学上彻底消灭消失/爆炸梯度。模型仍然可能因为门值设置不当、初始化不佳或序列太长而训练困难。

在此基础上，GRU（Gated Recurrent Unit）可以理解为更紧凑的门控变体：它把状态更新路径进一步简化，用更少的门与更少参数换来更轻量的实现。课程的重点仍放在 LSTM，但把 GRU 摆在同一脉络里，是为了强调“门控”才是核心思想，而不是某一个特定公式版本。

\subsection{本章小结}
LSTM/GRU 的本质不是把 RNN 变复杂，而是把状态更新从单次非线性变换改造成“受控的信息流”。一旦把长期记忆通道和即时输出通道拆开，模型就更有机会同时保留旧信息、吸收新信息并稳定训练。

\section{语言模型、图像描述与序列到序列}
% source_units: cu_07_11, cu_07_12, cu_07_13
本讲最后把 RNN 拉回到两个最典型的应用：语言建模（language modeling）与图像描述（image captioning）。

在语言模型里，训练目标通常是：给定前缀 $w_1, \dots, w_{t-1}$，预测下一个 token $w_t$。如果把前一个词的 embedding 作为输入，RNN 每一步都会更新状态并输出一个词表分布。训练时常用 teacher forcing，也就是把真实前缀喂给模型；推理时则改为把模型自己刚生成的 token 再送回去。

\begin{lstlisting}[caption={RNN decoder / language model 的典型 teacher forcing 训练流程}]
h = h0
loss = 0
for t in range(T):
    x_t = embed(tokens[t])
    h = rnn_cell(x_t, h)
    logits = W_hy @ h + b_y
    loss += cross_entropy(logits, tokens[t + 1])
\end{lstlisting}

图像描述则是在这个框架前面再接一个视觉编码器。最常见的思路是：
\begin{itemize}
\item 用 CNN 或其他视觉 backbone 把图像编码成向量或空间特征。
\item 用这个视觉表征初始化 decoder 的 hidden state，或者把它作为每步附加输入。
\item 让 RNN 逐词生成描述。
\end{itemize}

\begin{figure}[H]
\centering

\includegraphics[width=0.82\textwidth]{pdf_pages/page-90.png}
\caption{图像描述把视觉表征和词序列生成结合起来，是 RNN 在视觉任务中的经典应用。}
\end{figure}

这类方法的优点是结构直观、训练范式清晰，也天然适合 one-to-many 与 sequence-to-sequence 任务。但它也暴露出前面说过的瓶颈：所有历史上下文都被压进单条状态链中，越长的描述、越复杂的依赖、越难优化。

\begin{importantbox}{为什么下一讲会转向 attention}
当句子很长、输入很长，或 decoder 在某个时间步只需要关注输入中的局部区域时，单个 hidden state 往往成为瓶颈。attention 的价值，就是允许模型在每一步“直接回看”更丰富的上下文，而不必把一切都塞进同一个状态向量。
\end{importantbox}

\subsection{本章小结}
RNN 的应用逻辑非常统一：先维护状态，再按任务把状态读成分类、序列或文本。语言模型和图像描述说明了它的表达力，也把 hidden-state bottleneck 的问题暴露得最清楚。

\section{总结与延伸}
本讲把序列建模的基本问题讲得非常系统：为什么序列需要内部状态、vanilla RNN 如何通过参数共享处理变长输入、BPTT 为什么训练困难、LSTM/GRU 如何通过门控改善长依赖，以及这些结构如何落到语言模型与图像描述里。

从今天回看，RNN 的历史意义并不只在于它曾经主导序列任务，更在于它把几个后来一直延续的问题摆到了台面上：
\begin{itemize}
\item 状态应该如何压缩历史信息？
\item 长程依赖为什么难训练？
\item 计算图应该怎样在表达力、并行性与优化稳定性之间折中？
\end{itemize}

课程在结尾已经给出下一步方向：一方面，像 LSTM、GRU、modern state-space models（例如 Mamba）这样的结构继续探索更好的状态更新；另一方面，attention 则从根本上放宽了“所有历史都必须通过单一状态链传递”的限制。下一讲的核心，就是这条新路线如何出现并最终演化成 Transformer。

\end{document}
